\chapter{We All Descend from Migrants}
\section{A Stone Worked to an Edge}
Pick something up. It lies on my desk now: a dark-grey stone slightly smaller than an adult palm, shaped like an elongated drop, carefully worked on both faces, with an edge thin enough to cut paper.

It is a replica of a stone axe from Africa over a million years ago. Archaeologists call these Acheulean handaxes, after Saint-Acheul in France, where the type was first recognised. Making one requires more than casually striking a handy stone twice. The maker must first hold a plan in mind: the final teardrop, symmetrical curves, a thin continuous edge. Then another stone knocks off dozens of flakes, stroke by stroke. Each angle and force must be right; one error ruins the shape and wastes the effort.

Remember this: around 1.7 million years ago, a creature could already preserve a plan not visibly present and realise it through dozens of precise actions. That creature was our relative, or perhaps our ancestor.

More astonishingly, this handaxe form scarcely changed over the following million-plus years. Across Africa, Europe, and much of Asia, people made almost identical teardrops for over a million years; steam engines to smartphones took us little more than two centuries. Perhaps they were conservative; perhaps the design worked too well to need alteration. Either way, their makers were not beasts. Their world already contained tradition: one generation teaching a practice to another for over ten thousand generations.

This stone opens our chapter. Who made it, and how are they related to you and me, holding this book now?

\section[A Riverside Night Forty Millennia Ago]{A Riverside Night Forty Thousand Years Ago}
Enter a scene with me. Put aside your image of primitive humans; we will build this scene from nothing.

Time: an early-autumn night around forty thousand years ago. Place: a river in today's southwestern France. Shallow rock shelters line a chalk-white cliff; three fires burn on the open ground before them.

Woodsmoke mingles with roasting venison and a fishy smell: someone gutted a salmon earlier, leaving scales beside the fire to glitter. A dozen or more people sit around the flames. A wrinkled old woman sews hide clothing with a bone needle, making dense, even stitches with thread split from animal sinew. Two children toss pebbles and catch them on the backs of their hands. In the darkest corner, a young person strikes flint against flint, pausing every few strokes to squint at the edge and change hands to inspect it. This is the handaxe we met earlier, but intended as a grave offering for an elder who will be buried tomorrow.

When nobody speaks, you hear the river, cracking sparks, and a distant animal call from beyond the cliff top.

If you sit beside the fire, many things happen. They make room, offer roast meat, inspect your face and clothing, and speak an incomprehensible language with complete grammar, jokes, and a nickname for you. Next day, you might see someone enter the shelter's depths with a fat-burning lamp and paint a bison: charcoal sketch first, then blown ochre outlining it. The bison is better drawn than most people today could manage, because its artist has spent a lifetime watching bison more carefully than anyone.

Anatomically, these people are indistinguishable from us. Raise one of their newborns today and the child would grow, attend school, and use a phone like any other. The difference lies not in their bodies but in forty millennia of accumulation: no writing, metals, crops, or cities, all awaiting over a thousand later generations' gradual achievements.

Place this night beside that stone. The rest of our chapter answers their shared question.

\section[Where We Come From: Several Questions]{Where We Come From Is Not Just One Question}
Disassemble the question before rushing to answers.

"Where do we come from?" sounds singular, but contains at least three questions.

First, bodies: when and where did our species' bodily blueprint first appear? Our scientific name, Homo sapiens, means "wise human" in Latin.

Second, routes: how did those earliest people spread worldwide? Which paths led to East Asia and Australia, and who first entered the Americas, when?

Third, most persistent: what about me? Beijing, Sichuan, Henan---have my ancestors always lived here? Textbooks say Peking Man was the Chinese people's ancestor. Does that claim hold?

The third persists because it is personal. Your family genealogy may stop at the Qing; earlier paper records vanish. Yet your body holds an unbroken record whose answers, as we will see, differ from many peoples' cherished stories.

First establish a chapter-wide and book-wide rule: separate what happened, why, how participants understood it, and our evaluation. Fossils and genes can answer the first two. The wish that one's ancestors always lived here belongs to the latter two: real and understandable, but no substitute for evidence. Conversely, evidence must not mock that feeling.

Another agreement: many dates follow. Do not memorise precise numbers; scientists regularly revise them. Orders of magnitude suffice: our genus's story spans millions of years, our species' hundreds of thousands, the great African dispersal tens of thousands. Sequence matters more than numbers.

Lay those three questions on the table. Next we hear conflicting answers and learn to check evidence ourselves.

\section{Two Teams, Two Family Trees}
Twentieth-century palaeoanthropology roughly divided into two camps arguing over origins for decades. Hear both fully: their dispute itself teaches a lesson.

\textbf{Account 1: Out of Africa, newcomers prevail.}

Although members of Homo travelled between Africa and Eurasia much earlier, this camp locates all living humans' direct ancestors among Homo sapiens appearing in Africa only two or three hundred thousand years ago. Several branches dispersed on a large scale around sixty to seventy thousand years ago, following coasts and rivers worldwide. In Europe, Asia, and Australia, they replaced older populations, Europe's Neanderthals and Asian descendants of Homo erectus, who left few descendants. On this account, Zhoukoudian's Peking Man was no ancestor of ours, but a distant relative disappearing before our ancestors arrived.

\textbf{Account 2: Many regional blooms, one connected network.}

The multiregional hypothesis says Homo erectus had already left Africa over a million years earlier and spread across Eurasia. Regional populations stayed in contact and interbred within a vast network of gene flow. African, European, and East Asian populations retained some local differences while collectively evolving into modern Homo sapiens. Peking Man need not be your specific direct ancestor, but ancient East Asian populations genuinely contributed to today's East Asian ancestry.

Notice that I deliberately give the second account dignified treatment. Even a likely losing view deserves its strongest argument before testing. Multiregionalism was not mere amateur homesickness: it had fossil support. Certain East Asian facial and dental traits seemed continuous across hundreds of thousands of years, like inherited family features. An honest scientist could reasonably hypothesise regional population continuity from them.

Add something lest "out of Africa" sound like one event. History involved repeated departures. Over a million years ago Homo erectus left Africa, with fossils in Georgia's mountains, Java's riverbanks, and Zhoukoudian's caves. Homo sapiens also made multiple attempts: fossils over a hundred thousand years old in present-day Israel show earlier small departures, whose participants mostly left no continuous present-day descendants. Our main bodily lineage comes from the successful sixty-to-seventy-thousand-year dispersal. Some followed coasts eastward and sailed to Australia at least fifty thousand years ago, among the earliest evidence of human seafaring. Others hunted across icy terrain into Europe; others walked to the earth's ends. Migration is not some group's exceptional ordeal in this book, but ordinary life throughout three million years of Homo.

Then genetic evidence arrived.

From the late twentieth century, scientists directly compared living populations' genetic information. In the twenty-first, they even extracted surviving genetic material from ancient bone powder. This bodily genealogy broadly traces most living humans' genetic material to African Homo sapiens two or three hundred thousand years ago. Out-of-Africa got the main story right. Yet people outside Africa also carry a small component, roughly two per cent, from Neanderthals. Some Oceanian and Asian populations carry several per cent from another ancient population, Denisovans, identified only in 2010 through bone fragments and a genetic sequence from a Siberian cave.

Each camp thus won and lost something. Out-of-Africa won the main lineage: predominantly African ancestry. Multiregionalism recovered a corner: replacement was incomplete and encounters were not always deadly contests. Some encounters entered our bodies. This is no muddled compromise; evidence reshaped the question from who replaced whom to what happened when they met.

\section[Thinking Bridge: Evidence Has Jobs]{A Thinking Bridge: Assign Evidence Its Proper Work}
Grand ancestry questions invite a familiar trap: hear that thirty-thousand-year-old bones have been found somewhere and assume the entire family tree needs rewriting. A simple portable tool helps: \textbf{first identify the kind of evidence, then the question it can answer.}

Origins research relies chiefly on three kinds of evidence, like witnesses with different personalities.

\textbf{Witness 1: Fossils.} Bones and stones endure, revealing appearance, location, and approximate age. But fossils speak poorly: a skull cannot reveal language, singing, or relations with neighbours. Worse, the record is full of holes. Few creatures fossilise; fewer fossils are found. For a time, all material Denisovan evidence consisted of a few fragments. Fossils answer what and where well, but cannot alone readily answer what happened.

\textbf{Witness 2: Material remains.} Tools, hearths, paintings, beads, burials. These reveal behaviour: toolmaking implies planning and transmission; fire environmental modification; grave offerings a concept of afterlife; bison paintings the symbolic behaviour discussed later. But remains do not name their makers. A tool assemblage cannot identify who struck it.

\textbf{Witness 3: Genes.} Every living body is an archive; surviving ancient genetic material supplies original documents. Genes excel at kinship and ancestry flows, revealing Neanderthal ancestors where no fossil could. They fall almost silent about daily life, unable to report jokes around that forty-thousand-year-old fire.

See the method? Each witness has a remit and must not exceed it. A modern-looking fossil cannot establish direct ancestry; genetic checking is needed. Genetic resemblance cannot directly narrate a people's customs; material remains must contribute. The habit extends beyond archaeology. After news announces that research shows something, ask which research, and whether its evidence can actually answer the claimed question.

\section[Ancestors in Ancient Chinese Thought]{How Chinese People Imagined Ancestors Two Millennia Ago}
Read some primary material. Modern scientists are not alone in asking where we come from. Chinese thinkers over two thousand years ago lacked fossils and genes, but had answers. Late Warring States thinker Han Fei wrote in "The Five Vermin", Han Feizi:

\begin{quote}
In high antiquity people were few and beasts numerous, and people could not withstand beasts, insects, and snakes. A sage arose and built wooden nests against these harms. Delighted, the people made him ruler of the world and called him Youchao. People ate fruits, gourds, mussels, and clams; their raw, foul smells injured stomachs and caused much illness. A sage arose and drilled wood for fire to transform this raw food. Delighted, the people made him ruler and called him Suiren.
\end{quote}
The sense: outnumbered by animals, early people suffered until a sage taught wooden tree nests for protection. They acclaimed him leader, Youchao, the Nest-Maker. Raw produce and shellfish harmed stomachs and caused illness; another sage drilled fire and cooked food, earning leadership as Suiren, the Fire-Maker.

Read slowly; this is fascinating. Han Fei imagines antiquity not through divine creation but a plain idea: early life was hard, and nesting, fire, and cooking were invented step by step, not innate. Today that intuition seems strikingly close to reality: fire, shelter, and food processing were genuine historical thresholds. Without a single fossil, a thinker two millennia ago reasoned toward ancestors living differently and abilities accumulating afterward.

His era also marks the answer. Every skill comes from a sage, an acclaimed leader; exceptional individuals drive history. This reflects Warring States experience: rival states won through leadership and strategy, so antiquity becomes successive sagely kings' founding enterprises. People's accounts of origins carry their own age's colouring, as ours do today.

Clarify his purpose too: he was defending political principles, not writing prehistory. Changing times require changing government, not rigid old rules. The same essay says, "High antiquity competed in virtue; the middle age pursued strategy; today contends in strength." Earlier people compared virtue, later ones cunning, contemporaries power. His ancient past mirrors the present. Every age's origin story answers not merely what happened, but how we should live now.

\section[Their Disappearance: Three Accounts]{They Disappeared: The Same Bones, Three Accounts}
An earlier question remains: what happened when African Homo sapiens met Neanderthals in Europe and western Asia?

First establish evidence. Neanderthals, named after Germany's Neander Valley where they were first discovered, lived in Europe and western Asia for hundreds of thousands of years, much longer than sapiens there. Robust bodies, brains no smaller than ours, fine stone tools, fire, and care for injured companions characterised them. Iraq's Shanidar Cave yielded an individual who survived many years after severe injury, blindness in one eye, and an impaired arm, implying prolonged feeding and protection. Croatian eagle talons over a hundred thousand years old show wear plausibly from a necklace. Some buried their dead. Then around forty thousand years ago, shortly after sapiens entered Europe, their fossils and remains gradually disappear.

That disappearance has received at least three different accounts.

\textbf{Account 1: We defeated them.} Longer-range weapons, wider cooperation, stronger language, or denser trade-like exchanges helped sapiens outcompete Neanderthals for food and territory, perhaps through direct conflict. Timing makes this intuitive: sapiens arrives, Neanderthals leave. Limits are clear: almost no direct evidence of large-scale slaughter, and greater intelligence often mistakes result for cause. Because we survived, we must have been superior at everything: fluent reasoning, difficult to test.

\textbf{Account 2: Climate and population defeated them.} Europe's climate fluctuated severely, forests and ice repeatedly advancing and retreating. Neanderthals were always few and dispersed. Worse conditions could isolate small groups, triggering inbreeding and food-shortage cycles. Decline may have preceded sapiens, whose arrival was only the final straw. This avoids assuming superiority, but struggles to explain why people surviving hundreds of thousands of years of climatic swings failed this time.

\textbf{Account 3: They partly became us.} Genes show repeated interbreeding and descendants surviving today. Everyone outside Africa carries some Neanderthal ancestry. In 2018, scientists identified a Denisova Cave bone fragment from a teenage girl whose mother was Neanderthal and father Denisovan: two ancient populations' encounter embodied in one person's bones. Some disappearance was absorption. Strong new evidence supports this account, but not everything: why so little absorption, and why does only sapiens finally remain in the fossil record?

Each likely contains truth. Most researchers today favour a combination: small populations struggling through bad climate face newcomers' competition and mixing, then depart slowly and unevenly, enduring longer in some valleys than elsewhere. Notice the sentence form: several explanations each account for part of the facts. That is not evading a choice, but the most honest formulation present evidence permits.

\section[Further Challenge: Bush, Not Ladder]{A Deeper Challenge: Our Family Tree Is a Bush, Not a Ladder}
Now our strongest theoretical tool dismantles several widespread false pictures, including the caveman image we should have discarded at the start.

Describe it first: silhouettes from left to right, a hairy stooped ape gradually straightening and losing hair until a modern human stands proudly at the end. Innumerable books and advertisements imply three things: evolution is linear, moves from lower to higher, and ends with us as champions. All three are wrong.

Palaeontologists' real picture resembles a bush, not a ladder.

Consider facts. Since Homo, our small family, appeared in Africa over two million years ago, it has never consisted of one member alone. Multiple humans frequently shared the earth: Homo habilis and erectus in Africa over a million years ago; Neanderthal ancestors in Europe, erectus in Asia, sapiens ancestors in Africa hundreds of thousands ago. Even tens of thousands ago, sapiens, Neanderthals, Denisovans, and Indonesia's roughly metre-tall Homo floresiensis coexisted. The latter were no patients or mythical dwarfs but a distinct tool-using human species. One branch produces a species, another a different one; some flourish, some wither. More than a dozen species appeared over time. We are not a lone shoot reaching the top, but the bush's last branch surviving today.

This corrects three everyday mistakes.

Mistake one: humans evolved from monkeys. No: today's monkeys and apes are our cousins. Our common ancestor was neither a present-day monkey nor chimpanzee, just as sharing a grandfather with your cousin does not mean evolving from your cousin.

Mistake two: erectus became sapiens in sequential stages. Probably wrong too. Erectus remained in East Asia hundreds of thousands of years ago, while our direct ancestors were African. They resemble two branches more than consecutive relay runners.

Mistake three: later arrivals are more advanced. Evolution has no direction called advanced, only adaptation to current conditions. Mammals rose after dinosaurs disappeared not because they were better, but because the environment changed the question. Imagining evolution as levelling up mistakes outcome for purpose. Evolution has no purpose, only survival.

The savage-caveman stereotype follows ladder thinking: if we occupy the top, everything earlier must be rougher, dimmer, more brutal. Recall the evidence: symmetrical handaxe makers calculated many striking angles; Shanidar Neanderthals sustained a badly injured companion; forty-thousand-year-old artists painted bison better than most of us. Early twentieth-century scholars reconstructed a French Neanderthal elder as a hunched, bent-kneed, repulsive savage, influencing generations of textbooks. Re-examination found severe joint disease in an already elderly individual. Representing a species with a frail old patient resembles representing humanity through hospital patients. The stereotype was not excavated from stone: we bent fossils to fit our ladder.

The theory also has a useful by-product. You may know mitochondrial Eve: genetic tracing converges all living people's maternal lineages on an African woman over a hundred thousand years ago. Does that mean only one woman remained alive? No. Thousands of others lived then, but their strictly maternal lines eventually ended through only sons or no descendants. Eve's happened to remain unbroken. She is a genealogical convergence point, not the sole survivor or first woman. A prominent tree node never means only one creature then existed worldwide. That is ladder thinking in another costume.

\section{From Zhoukoudian to Your Saliva}
This ancient question now comes as close as a tube of saliva.

Consumer genetic tests have become popular: mail saliva and receive a report on ancestral origins weeks later. Feelings can be complicated. Supposedly rooted families discover components from several directions; attempts to prove pure blood reveal unexpected ancestry. Our tools suggest two cautions. First, tests measure resemblance to reference populations: shared ancestry with people in certain regions today, not a pure family line. Second, pure blood is almost a genealogical non-concept. Trace upstream through the bush and everyone's ancestors spread into a vast network, each node occupied by a living person who migrated, intermarried, and told children stories.

A larger echo comes from national narratives. Nearly every country has cherished a story of uninterrupted local residence and shared blood. China has called Peking Man Chinese people's ancestor; Europe has romantically adopted Neanderthals as indigenous forebears, or darkly portrayed newcomers as contaminating invaders. Evidence challenges both. Against immobility: every population moved. Erectus travelled over a million years ago; sapiens left Africa tens of thousands ago, entered Australia over subsequent millennia, and crossed the then-land-connected Bering Strait into the Americas. Your ancestors, like everyone's, were travelling. Against contamination: encounter and mixture are ordinary history, not exceptions. Your body may carry an encounter over forty thousand years old.

Recall our fourfold distinction. Those conclusions interpret evidence. How to treat present-day migrants or regard one's nation is evaluation, where evidence cannot decide for you. Science can expose pure-blood misconceptions but cannot directly determine immigration policy. Ethics and common sense must help separate what happened from what should happen. That is humanistic training's proper work.

Bring the camera into your classroom. Has someone asked where your family comes from? Tone changes the question entirely. Fellow locals seek connection; city pupils questioning a new classmate sometimes classify; online arguments often preface expulsion: "People like you should go back where you came from." A species migrating for hundreds of thousands of years still divides us from them through origins. That fact may deserve longer attention than any fossil.

\section[Your Turn: How Long Makes a Local?]{Your Turn: How Long Must You Stay to Count as Local?}
Return to the first handaxe.

If its maker could know her work would lie in a museum over a million years later, she probably would not imagine descendants arguing so bitterly over who belongs where.

End with a question having no standard answer:

\textbf{How long must a group inhabit a place before it counts as local?}

Answer with reasons, first weighing these considerations again.

If arrival order decides, everyone outside Africa descends from arrivals within tens of thousands of years. Indigenous Americans' ancestors crossed ten or twenty thousand years ago. Even African populations never stopped moving. Almost any chosen cutoff year is arbitrary.

If cultural inheritance decides, culture itself changes along the journey. Languages, recipes, gods change; no people preserves one cultural package untouched for millennia.

If you abandon the boundary and say everyone migrates and belongs everywhere, that sounds generous but erases something real: long residence makes a place irreplaceable, holding graves, dialect, memory. "Everyone came from elsewhere anyway" sounds weightless to someone just expelled from home.

Should local identity follow years of residence, ancestry, knowledge and love of land, or no one's unilateral decision? When "I am local" welcomes others, it is kindness. What is it when used to expel them?

No standard answer exists. But each word you use---we, they, ancestors, home---did not fall from the sky. Each is a fossil of hundreds of thousands of years of migration surviving in your mouth. From today, you have eyes to recognise them.
